% ID: SEL-D0047-OB-E02
% ESTADO: validada. S048.
\textbf{a)} En el vértice se cumple $f'(-1)=2+p=0$, luego $p=-2$. El coeficiente cuadrático negativo garantiza un máximo. Como $f(-4)=-16+8+q=-5$, resulta $q=3$. El valor máximo es $f(-1)=-1+2+3=\boxed{4}$.

\textbf{b)} Para $p=2$ y $q=-1$ queda $f(x)=-(x-1)^2$. Es una parábola cóncava hacia abajo, con vértice $(1,0)$, eje $x=1$ y corte $(0,-1)$ con el eje vertical. En $x=-2$ se tiene $f(-2)=-9$ y $f'(-2)=6$. Así, $y+9=6(x+2)$, es decir, $\boxed{y=6x+3}$.
\begin{center}
\setlength{\unitlength}{0.9mm}
\begin{picture}(94,55)(-35,-45)
\qbezier(-34,0)(10,0)(54,0)
\qbezier(0,-43)(0,-18)(0,7)
\put(55,1){$x$}
\put(2,7){$y$}
\qbezier(-32,-36)(8,36)(48,-36)
\put(6,2){\scriptsize $(1,0)$}
\put(-34,-39){\scriptsize $(-2,-9)$}
\end{picture}
\end{center}
